% 解析論文の主張・結果・考察を，予備知識の無い大学 1 年生向けに説明する資料（2026-11-13）
% 作り方：cd DAI/docs && lualatex 説明_主張_大学1年生向け_2026-11-13.tex （2 回）
\documentclass[a4paper,11pt]{ltjsarticle}
\usepackage[margin=22mm]{geometry}
\usepackage{booktabs}
\usepackage{float}
\usepackage{array}
\usepackage{xcolor}
\usepackage{tikz}
\usetikzlibrary{arrows.meta,positioning,fit,calc}
\usepackage{pgfplots}
\pgfplotsset{compat=1.18}
\usepackage{tcolorbox}
\usepackage{enumitem}
\usepackage[hidelinks]{hyperref}
\setlist{itemsep=2pt,topsep=4pt}

\definecolor{cA}{HTML}{2A6FB0}  % 声
\definecolor{cV}{HTML}{D2691E}  % 顔
\definecolor{cT}{HTML}{3A8F3A}  % 言葉
\definecolor{cF}{HTML}{7B4FA0}  % 融合
\definecolor{gr}{HTML}{888888}

\newtcolorbox{yoten}[1]{colback=black!3,colframe=black!55,boxrule=0.6pt,arc=1mm,title=#1,fonttitle=\bfseries}

\title{新しいコーパスに感情ごと 5 本しかラベルが無いとき，\\疑似ラベルで学び直すと UAR が上がるのはどの分類器か\\[4pt]
\large 解析論文の主張・実験の結果・考察の説明}
\author{MER-DA（DAI の作業場）}
\date{2026 年 11 月 13 日}

\begin{document}
\maketitle

\begin{yoten}{この資料の要点}
\begin{itemize}
  \item 新しい感情のデータ（ターゲットのコーパス）に，感情ごとに 5 本しか正解のラベルが無い場面を調べた．調べた方法（以下「提案」）は，声・顔・言葉の 3 つの分類器の予測を掛け合わせて，ラベルの無い発話に仮の正解（疑似ラベル）を付け，その疑似ラベルで分類器を学び直す．
  \item 4 つのコーパスで試すと，UAR（感情ごとの正解率の平均．でたらめなら 25 点）が 4 コーパスすべてで上がったのは声の分類器だけだった．上がり幅は +2.4〜+4.1 点である．ただし IEMOCAP の +2.4 点は区間が 0 を含む．
  \item 論文の主張 1：学び直して UAR が上がるのは，「余地がある」「疑似ラベルが今の分類器より当たる」「疑似ラベルの判断を自分の特徴で真似できる」の 3 つの条件がそろうときである．4 コーパスすべてで 3 つがそろったのは声の分類器だった．
  \item 論文の主張 2：声の特徴は，言葉の分類器の判断を，顔の特徴より 15〜17 点よく真似できた．感情と話者が同じことで説明できる分を差し引いても，10〜17 点の差が残った．
  \item 弱いところ：3 つの条件は結果を見た後に立てた説明であり，同じ 4 コーパスと 1 組のエンコーダでしか確かめていない．
\end{itemize}
\end{yoten}

\tableofcontents
\newpage

%======================================================================
\section{何を調べたか}

\subsection{感情認識と，成績の測り方}

この研究は，人が話している短い動画（\textbf{発話}．数秒の 1 区切り）を見て，話し手の感情を当てるプログラムを扱う．感情は Neutral（平静），Happy（喜び），Sad（悲しみ），Angry（怒り）の 4 つに限る．

発話からは 3 種類の情報が取れる．この情報の種類を\textbf{モダリティ}と呼ぶ．
\begin{itemize}
  \item \textcolor{cA}{\textbf{声}}：音声の波形．声の高さ，強さ，話し方，そして話している言葉の音も含む．
  \item \textcolor{cV}{\textbf{顔}}：映像に写る顔の表情．
  \item \textcolor{cT}{\textbf{言葉}}：話した内容の文字．人が書き起こしたもの（AFEW だけは音声認識の Whisper が書き起こしたもの）．
\end{itemize}
3 つをまとめて使う分類器を\textbf{融合の分類器}と呼ぶ．

成績は\textbf{UAR}で測る．UAR は，4 つの感情それぞれの正解率を出し，その 4 つを平均した値である．たとえば Neutral の発話の 80\%，Happy の 60\%，Sad の 40\%，Angry の 20\% を当てたら，UAR は (80 + 60 + 40 + 20) ÷ 4 = 50 点になる．でたらめに答えると，どの感情も 25\% しか当たらないので UAR は 25 点になる．ふつうの正解率ではなく UAR を使うのは，データの中で Neutral の発話だけが多い場合に，「全部 Neutral と答える」だけで高い点が出るのを防ぐためである．

\subsection{困っている場面}

感情のラベル（人が付けた正解）が付いた発話の集まりを\textbf{コーパス}と呼ぶ．この研究では公開されている 4 つのコーパスを使った．
\begin{itemize}
  \item \textbf{AFEW}：映画の場面の切り出し（約 800 発話）．
  \item \textbf{IEMOCAP}：俳優が実験室で演じた会話．この研究では即興の部分を使う．
  \item \textbf{MELD}：テレビドラマ「Friends」の台詞（約 1 万 1 千発話）．
  \item \textbf{MOSEI}：YouTube の動画で人が意見を話す場面（約 1 万 9 千発話）．
\end{itemize}

新しい場面（新しいコーパス）で感情認識を使いたいとき，正解のラベルを大量に付けるには人手と時間がかかる．そこで，次の場面を想定した．
\begin{itemize}
  \item 新しいコーパス（\textbf{ターゲット}）には，感情ごとに 5 本，合わせて 20 本の発話にしかラベルが無い．この 20 本を「K 本」と呼ぶ（K = 5）．
  \item ターゲットには，ラベルの無い発話はたくさんある．
  \item ほかのコーパス（\textbf{ソース}）には，ラベルが全部付いている．
\end{itemize}
4 つのコーパスのうち 1 つをターゲット，残りの 3 つをソースにする．これを 4 通り繰り返す．この組み方を\textbf{LOCO 設定}（leave-one-corpus-out．1 つのコーパスを抜く設定）と呼ぶ．

\subsection{使った道具}

発話をそのまま分類器に入れることはせず，まず\textbf{エンコーダ}で数字の並び（\textbf{特徴}）に変える．エンコーダは，大量のデータで既に学習してあるニューラルネットである．この研究では，エンコーダの中身を一切変えずに使う（\textbf{凍結}）．
\begin{itemize}
  \item 声：WavLM-large の最後の層（1 発話を 1024 個の数に変える）．
  \item 顔：ResNet-50．表情の写真のデータ（AffectNet）で表情を当てるように学習したもの（512 個の数）．
  \item 言葉：RoBERTa-base の最後の層（768 個の数）．
\end{itemize}
特徴の上には\textbf{線形の分類器}を置く．線形の分類器は，特徴の各数に感情ごとの重みを掛けて足し合わせ，その合計（点数）が一番大きい感情を答える．学習とは，この重みを決めることである．

\subsection{実験の組み方}

\begin{itemize}
  \item \textbf{話者を分ける}：同じ人の発話が，学習に使う側と成績を測る側（\textbf{テスト}）の両方に入らないようにした．同じ人の声や顔を覚えただけで点が上がるのを防ぐためである．MOSEI には話者の情報が無いので，YouTube のチャンネル（無ければ動画）で分けた．
  \item \textbf{繰り返し}：話者を 5 つの組に分け，1 組ずつテストにする．さらに K 本の選び方を 10 通り変える．合わせて 50 回の結果から平均と区間を出した．
  \item \textbf{区間}：この資料の [ ] の中の 2 つの数は 95\% 区間である．50 回の結果のばらつきから作った幅で，「本当の差はおそらくこの範囲にある」と見込む範囲を表す．区間が 0 を含まなければ「上がった（下がった）」と言い，0 を含めば「判定できない」と言う．「判定できない」は「差が無い」とは違う．
  \item \textbf{設定の選び方}：方法の設定（重みの強さなど）は，評価に使うコーパスのテストの正解を見ずに決めた．ソースの 1 つを仮のターゲットにした試し（\textbf{仮のターゲットの回}）で決めて固定し，そのあとで評価の回を 1 回だけ流した．
\end{itemize}

%======================================================================
\section{調べた方法（提案）}

\begin{figure}[H]
\centering
\begin{tikzpicture}[
  box/.style={draw,rounded corners=2pt,align=center,minimum height=9mm,inner sep=3pt,font=\small},
  arr/.style={-{Stealth[length=2.2mm]},thick}]
  \node[box,fill=black!4] (src) {ソースの 3 コーパス\\（ラベルは全部ある）};
  \node[box,fill=black!4,right=9mm of src] (start) {ソースで学んだ\\分類器の平均\\（出発点）};
  \node[box,fill=black!4,below=10mm of src] (k) {ターゲットの K 本\\（感情ごと 5 本）};
  \node[box,right=9mm of k] (base) {K 本だけで学び直した\\声・顔・言葉の分類器};
  \node[box,fill=yellow!15,right=9mm of base] (pl) {3 つの予測を掛け合わせた\\\textbf{疑似ラベル}};
  \node[box,fill=black!4] (un) at (pl |- start) {ターゲットの\\ラベル無しの発話};
  \node[box,fill=blue!6,below=10mm of pl] (re) {K 本 ＋ 疑似ラベルで\\各分類器を学び直す（\textbf{提案}）};
  \draw[arr] (src) -- (start);
  \draw[arr] (start) -- (base);
  \draw[arr] (k) -- (base);
  \draw[arr] (base) -- (pl);
  \draw[arr] (un) -- node[right,font=\scriptsize]{予測を付ける} (pl);
  \draw[arr] (pl) -- (re);
  \draw[arr] (k.south) |- (re.west);
\end{tikzpicture}
\caption{提案の流れ．比べる相手の「K 本だけ」は，出発点から K 本だけで学び直した分類器である．テストの発話は，この図のどこにも入らない．}
\label{fig:flow}
\end{figure}

提案の手順は次のとおりである（図\ref{fig:flow}）．
\begin{enumerate}
  \item ソースの 3 コーパスで声・顔・言葉の分類器をそれぞれ学び，その重みを平均する．これを出発点にする．
  \item 出発点から，ターゲットの K 本（20 本）だけで分類器を学び直す．これが比べる相手の「\textbf{K 本だけ}」の分類器である．
  \item ラベルの無い発話 1 本ごとに，3 つの分類器が出した「各感情である確率」を掛け合わせ，合計が 1 になるように割り直す．これが疑似ラベルである．
  \item K 本と疑似ラベルの両方を使って，各分類器を出発点から学び直す．これが提案の分類器である．
\end{enumerate}

\paragraph{疑似ラベルの作り方の数値例}
ある発話について，3 つの分類器が次の確率を出したとする．
\begin{center}
\begin{tabular}{lcccc}
\toprule
 & Neutral & Happy & Sad & Angry \\
\midrule
\textcolor{cA}{声の分類器} & 0.4 & 0.3 & 0.2 & 0.1 \\
\textcolor{cV}{顔の分類器} & 0.3 & 0.4 & 0.2 & 0.1 \\
\textcolor{cT}{言葉の分類器} & 0.5 & 0.2 & 0.2 & 0.1 \\
\midrule
3 つの積 & 0.060 & 0.024 & 0.008 & 0.001 \\
割り直した値（疑似ラベル） & \textbf{0.645} & 0.258 & 0.086 & 0.011 \\
\bottomrule
\end{tabular}
\end{center}
積の合計は 0.093 なので，各値を 0.093 で割ると疑似ラベルになる．3 つの分類器がそろって高い確率を付けた感情（この例では Neutral）ほど，疑似ラベルの確率が高くなる．
（実際の手順では，3 つの分類器を融合の分類器の予測で一度学び直してから掛け合わせ，掛け合わせるときの各分類器の効き具合を仮のターゲットの回で選んだ．流れは上と同じである．）

\paragraph{提案と普通の蒸留の関係}
ほかの分類器の予測を目標にして分類器を学ぶ方法は，一般に\textbf{蒸留}と呼ばれる．提案は，同じ疑似ラベルを渡す「温度つき蒸留」とほぼ同じ UAR になった（3 節）．このため論文は「提案が既存の方法より良い」とは主張しない．論文の主張は，「疑似ラベルで学び直したときに，どの分類器でなぜ UAR が上がるのか」の説明である．

%======================================================================
\section{結果 1：学び直して UAR はどれだけ上がったか}

\begin{table}[H]
\centering
\caption{UAR の変化（提案 − K 本だけ．K = 5．評価の回の 50 回）．単位は点．[ ] は 95\% 区間．\textbf{太字}は区間が 0 を含まない（上がった）もの．}
\label{tab:gain}
\small
\begin{tabular}{lcccc}
\toprule
学び直した分類器 & AFEW & IEMOCAP & MELD & MOSEI \\
\midrule
\textcolor{cA}{声} & \textbf{+4.1} [2.6, 5.7] & +2.4 [\textminus{}0.4, 5.3] & \textbf{+3.1} [1.2, 5.1] & \textbf{+3.2} [2.4, 4.0] \\
\textcolor{cV}{顔} & +0.1 [\textminus{}3.0, 3.2] & \textminus{}0.8 [\textminus{}4.7, 3.1] & +0.8 [\textminus{}0.9, 2.6] & +0.2 [\textminus{}1.1, 1.6] \\
\textcolor{cT}{言葉} & +1.6 [\textminus{}1.5, 4.8] & \textbf{+7.3} [4.4, 10.3] & \textminus{}1.2 [\textminus{}4.4, 2.1] & \textminus{}0.0 [\textminus{}0.9, 0.8] \\
\textcolor{cF}{融合} & +1.0 [\textminus{}0.7, 2.7] & +0.2 [\textminus{}4.5, 4.8] & \textbf{+2.5} [1.7, 3.4] & \textminus{}0.4 [\textminus{}0.9, 0.2] \\
\midrule
\multicolumn{5}{l}{\small 参考：声の分類器の UAR（K 本だけ → 提案）} \\
 & 50.1 → 54.2 & 67.1 → 69.5 & 41.5 → 44.6 & 34.9 → 38.1 \\
\bottomrule
\end{tabular}
\end{table}

\begin{figure}[H]
\centering
\begin{tikzpicture}
\begin{axis}[
  width=0.95\textwidth, height=6.2cm,
  ybar, bar width=7pt,
  ymin=-6, ymax=11,
  ylabel={UAR の変化（点）},
  symbolic x coords={AFEW,IEMOCAP,MELD,MOSEI},
  xtick=data, enlarge x limits=0.18,
  legend style={at={(0.5,1.03)},anchor=south,legend columns=4,font=\small,draw=none},
  extra y ticks={0}, extra y tick style={grid=major,grid style={black!60}},
  ymajorgrids, grid style={black!10},
  tick label style={font=\small}, label style={font=\small},
  error bars/y dir=both, error bars/y explicit,
]
\addplot[fill=cA,draw=cA,error bars/error bar style={black,thin}] coordinates {
 (AFEW,4.150) += (0,1.589) -= (0,1.593) (IEMOCAP,2.413) += (0,2.861) -= (0,2.860) (MELD,3.118) += (0,1.938) -= (0,1.939) (MOSEI,3.187) += (0,0.803) -= (0,0.803)};
\addplot[fill=cV,draw=cV,error bars/error bar style={black,thin}] coordinates {
 (AFEW,0.092) += (0,3.081) -= (0,3.082) (IEMOCAP,-0.844) += (0,3.901) -= (0,3.901) (MELD,0.850) += (0,1.764) -= (0,1.765) (MOSEI,0.211) += (0,1.349) -= (0,1.348)};
\addplot[fill=cT,draw=cT,error bars/error bar style={black,thin}] coordinates {
 (AFEW,1.647) += (0,3.147) -= (0,3.147) (IEMOCAP,7.343) += (0,2.969) -= (0,2.969) (MELD,-1.178) += (0,3.259) -= (0,3.259) (MOSEI,-0.025) += (0,0.870) -= (0,0.869)};
\addplot[fill=cF,draw=cF,error bars/error bar style={black,thin}] coordinates {
 (AFEW,1.031) += (0,1.684) -= (0,1.683) (IEMOCAP,0.157) += (0,4.689) -= (0,4.689) (MELD,2.534) += (0,0.884) -= (0,0.884) (MOSEI,-0.354) += (0,0.544) -= (0,0.544)};
\legend{声,顔,言葉,融合}
\end{axis}
\end{tikzpicture}
\caption{表\ref{tab:gain}の値の図．黒い縦の線は 95\% 区間．}
\label{fig:gain}
\end{figure}

表\ref{tab:gain}と図\ref{fig:gain}に，提案の分類器の UAR から K 本だけの分類器の UAR を引いた値（\textbf{gain}）を示す．
\begin{itemize}
  \item \textbf{声の分類器}は 4 コーパスすべてで UAR が上がった（+2.4〜+4.1 点）．区間が 0 を含まないのは AFEW，MELD，MOSEI の 3 つで，IEMOCAP の区間は \textminus{}0.4〜+5.3 点と 0 を含む．
  \item \textbf{顔の分類器}は，4 コーパスすべてで区間が 0 を含んだ．
  \item \textbf{言葉の分類器}は，IEMOCAP だけで +7.3 点上がった．ほかの 3 つは区間が 0 を含んだ．
  \item \textbf{融合の分類器}は，MELD だけで +2.5 点上がった．ほかの 3 つは区間が 0 を含んだ．
  \item 4 コーパスすべてで値が正だったのは声の分類器だけである．
\end{itemize}

\paragraph{普通の蒸留との比べ}
同じ疑似ラベルを温度つき蒸留に渡した分類器と提案の分類器を比べると，比べた 48 の条件（4 コーパス × 4 種類の分類器 × K が 0，5，25 の 3 通り）のうち 40 で区間が 0 を含んだ．たとえば声の分類器（K = 5）の UAR は，AFEW で提案 54.2 点と蒸留 53.6 点，MOSEI で 38.1 点と 38.2 点だった．提案の値は普通の蒸留の値とほとんど変わらない．

%======================================================================
\section{主張 1：学び直して UAR が上がるための 3 つの条件}

\subsection{3 つの条件}

表\ref{tab:gain}を見ると，「なぜ声だけが 4 コーパスで上がり，顔は上がらないのか」が問いになる．この問いに対する論文の答えが次の 3 つの条件である．疑似ラベルで学び直して UAR が上がるのは，3 つがすべて成り立つときである．

\begin{enumerate}
  \item \textbf{余地がある．}\\
  ターゲットの全部の発話にラベルを付けて学べば，今の分類器より UAR が上がる．言いかえると，特徴の中に，まだ使い切っていない感情の情報が残っている．\\
  例：K 本だけの分類器が 50 点で，全部のラベルで学んでも 51 点にしかならないなら，どんな方法で学び直しても 1 点ほどしか上がらない．\\
  測り方：ラベルの無い発話の全ラベル（評価のときだけ使える）で学んだ分類器の UAR から，K 本だけの分類器の UAR を引く．

  \item \textbf{疑似ラベルが今の分類器より当たる．}\\
  学び直しでは，分類器は疑似ラベルに近い答えを出すように変わる．疑似ラベルの UAR が今の分類器より低ければ，分類器は悪くなる．\\
  例：顔の分類器がもともと 71 点で，疑似ラベルが 73 点なら，近づいても 2 点ほどしか上がらない．\\
  測り方：ラベルの無い発話の上で，疑似ラベルの UAR から K 本だけの分類器の UAR を引く（\textbf{当たる幅}）．

  \item \textbf{疑似ラベルの判断を，自分の特徴で真似できる．}\\
  線形の分類器は，自分の特徴の数に重みを掛けて足すことしかできない．疑似ラベルが「言葉の内容から Happy」と判断した発話があっても，顔の特徴にその手がかりが無ければ，顔の分類器はその判断を学べない．\\
  測り方：自分の特徴だけで疑似ラベルを当てるように学んだ線形の分類器を作り，その UAR から K 本だけの分類器の UAR を引く（\textbf{真似できる分}）．真似できる分は「疑似ラベルを完全に真似できたら，どこまで上がるか」の目安になる．
\end{enumerate}

\subsection{4 コーパスでの値}

\begin{table}[H]
\centering
\small
\caption{3 つの条件の値（K = 5．評価の回）．単位は点．[ ] は 95\% 区間．「余地」には区間を付けていない．最後の列は，下調べ（仮のターゲットの回）の値から，欠けていると判断した条件の番号．}
\label{tab:cond}
\setlength{\tabcolsep}{3pt}
\begin{tabular}{llcccccc}
\toprule
分類器 & コーパス & 条件 1 & \multicolumn{2}{c}{条件 2} & 条件 3 & gain & 欠けた \\
\cmidrule(lr){4-5}
 & & 余地 & K 本だけ → 疑似ラベル & 当たる幅 & 真似できる分 & （表\ref{tab:gain}） & 条件 \\
\midrule
\textcolor{cA}{声} & AFEW    & 6.1  & 51.0 → 72.7 & +21.7 [18.1, 25.3] & 3.7 [2.0, 5.4] & +4.1 & なし \\
                   & IEMOCAP & 10.1 & 68.1 → 75.4 & +7.2 [1.5, 13.0]  & 0.4 [\textminus{}3.8, 4.5] & +2.4 & なし \\
                   & MELD    & 6.5  & 42.3 → 50.6 & +8.4 [6.9, 9.8]   & 2.0 [0.9, 3.1] & +3.1 & なし \\
                   & MOSEI   & 11.4 & 34.6 → 41.3 & +6.6 [5.6, 7.7]   & 2.6 [1.8, 3.3] & +3.2 & なし \\
\midrule
\textcolor{cV}{顔} & AFEW    & 4.9  & 71.1 → 72.7 & +1.7 [\textminus{}4.0, 7.3]  & 1.3 [\textminus{}0.8, 3.4] & +0.1 & 2 \\
                   & IEMOCAP & 4.7  & 57.8 → 75.4 & +17.6 [13.0, 22.2] & \textminus{}1.1 [\textminus{}3.2, 1.0] & \textminus{}0.8 & 1・3 \\
                   & MELD    & 1.8  & 38.0 → 50.6 & +12.6 [11.0, 14.3] & 0.5 [\textminus{}0.7, 1.7] & +0.8 & 1 \\
                   & MOSEI   & 8.7  & 36.4 → 41.3 & +4.9 [3.3, 6.4]   & 0.3 [\textminus{}0.8, 1.5] & +0.2 & 3 \\
\midrule
\textcolor{cT}{言葉} & AFEW    & 3.7  & 45.4 → 72.7 & +27.3 [22.5, 32.1] & 2.8 [1.3, 4.4] & +1.6 & 3 \\
                     & IEMOCAP & 11.5 & 50.3 → 75.4 & +25.0 [22.4, 27.7] & 6.4 [2.8, 10.0] & +7.3 & なし \\
                     & MELD    & 7.9  & 50.3 → 50.6 & +0.3 [\textminus{}1.0, 1.7]  & \textminus{}1.4 [\textminus{}2.3, \textminus{}0.6] & \textminus{}1.2 & 2 \\
                     & MOSEI   & 7.9  & 33.4 → 41.3 & +7.9 [6.6, 9.2]   & \textminus{}0.2 [\textminus{}1.0, 0.6] & \textminus{}0.0 & 3 \\
\bottomrule
\end{tabular}
\end{table}

表\ref{tab:cond}に，声・顔・言葉の分類器について 3 つの条件の値を示す．「K 本だけ → 疑似ラベル」の列は，ラベルの無い発話の上での，K 本だけの分類器の UAR と疑似ラベルの UAR である．疑似ラベルは 3 つの分類器の予測から作るので，同じコーパスでは声・顔・言葉で同じ値になる．

\paragraph{声の分類器}
4 コーパスとも 3 つの条件がそろった．余地は 6.1〜11.4 点あり，疑似ラベルは声の分類器より 6.6〜21.7 点よく当たり，真似できる分の値は 0.4〜3.7 点で正だった（IEMOCAP は区間が 0 を含む．4.3 節）．gain（+2.4〜+4.1 点）は，真似できる分と同じくらいの大きさである．

\paragraph{顔の分類器}
コーパスごとに欠けた条件が違った．
\begin{itemize}
  \item AFEW：顔の分類器がもともと 71.1 点で，4 つの分類器の中で最も強い．疑似ラベル（72.7 点）との差は 1.7 点しかなく，区間は 0 を含む（条件 2 が欠ける）．
  \item IEMOCAP：疑似ラベルは顔の分類器より 17.6 点よく当たるが，真似できる分は \textminus{}1.1 点だった（条件 3 が欠ける）．下調べでは，全ラベルで学んでも分類器の重みの向きがほとんど良くならなかった（向きを直して得られる分は 0.4 点）ので，条件 1 も欠けると判断した．
  \item MELD：余地が 1.8 点しかない（条件 1 が欠ける）．
  \item MOSEI：真似できる分が 0.3 点しかない（条件 3 が欠ける）．
\end{itemize}

\paragraph{言葉の分類器}
\begin{itemize}
  \item MELD：言葉の分類器（50.3 点）が疑似ラベル（50.6 点）とほぼ同じだけ当たる（条件 2 が欠ける）．
  \item MOSEI：真似できる分が \textminus{}0.2 点である（条件 3 が欠ける）．
  \item AFEW：下調べでは真似できる分が 1.3 点と小さく，条件 3 が欠けると判断した．しかし評価の回で測り直すと 2.8 点 [1.3, 4.4] あり，この判断と合わない．
  \item IEMOCAP：3 つの条件がそろい，gain は +7.3 点だった．ただし IEMOCAP の即興の会話は，俳優の 2 つの組が，同じ 8 つの筋書きを演じたものである．このため，言葉の分類器は筋書きの話題から感情を当てている可能性があり，この上がり幅を疑似ラベルの効果だけに帰すことはできない．
\end{itemize}

\paragraph{融合の分類器}
融合の分類器は，初めから声・顔・言葉の 3 つの特徴を見ている．疑似ラベルも同じ 3 つのモダリティの予測から作るので，疑似ラベルは融合の分類器とほぼ同じだけしか当たらない．ラベルの無い発話の上で，疑似ラベルの UAR から融合の分類器（K 本だけ）の UAR を引いた差は +0.1〜+2.5 点で，2 つの判断は発話の 80〜89\% で一致した（条件 2 が欠ける．MOSEI だけは 50 回のうち 39 回の値）．例外は MELD で，下調べでは 3 つの条件がそろい，gain も +2.5 点あった．

\subsection{この表を読むときの注意}
\begin{itemize}
  \item 「欠けた条件」の列は，下調べの値を見て付けた．評価の回で測り直すと，声の真似できる分は 0.4〜3.7 点で，下調べの値（2.9〜4.5 点）より小さかった．特に IEMOCAP の声は 0.4 点で，区間は \textminus{}3.8〜+4.5 点と 0 を含む．このため IEMOCAP の声で条件 3 が満たされたとは，評価の回の値だけからは言えない．
  \item 3 つの条件は，表\ref{tab:gain}の結果を見た後に立てた説明である．計算の前に立てた予言を確かめたものではない．
  \item 4 コーパスすべてで 3 つの条件がそろったのは声の分類器だけである．言葉の分類器は IEMOCAP で，融合の分類器は MELD で，それぞれ 1 コーパスだけそろった．
\end{itemize}

%======================================================================
\section{主張 2：声の特徴は，言葉の分類器の判断を顔の特徴より良く真似できる}

\subsection{なぜこれを確かめたか}

主張 1 では，声の分類器だけが 4 コーパスで条件 3（疑似ラベルの判断を真似できる）を満たした．その理由として，研究者は次の説明を立てた．
\begin{quote}
音声には声色と話している言葉の両方が入っている．だから言葉の分類器から来た判断の一部は，声の特徴で真似できる．
\end{quote}
疑似ラベルは声・顔・言葉の 3 つの分類器の予測の積なので，言葉の分類器の判断も入っている．この説明が正しければ，「声の特徴は，言葉の分類器の判断を，顔の特徴より良く真似できる」はずである．この予言を計算の前に書いて固定し，それから確かめた．

\subsection{測り方}

\begin{enumerate}
  \item K 本だけで学んだ言葉の分類器に，ラベルの無い発話それぞれの感情を答えさせる．この答えを「言葉の分類器の判断」と呼ぶ．
  \item ラベルの無い発話を 2 つに分ける．片方の半分で，声の特徴から言葉の分類器の判断を当てる線形の分類器を学ぶ．
  \item もう片方の半分で，その分類器の答えが言葉の分類器の判断と一致するかを数える．一致は UAR と同じく感情ごとの一致の割合の平均で測る（\textbf{一致度}．でたらめなら 25 点）．
  \item 声の特徴の代わりに顔の特徴を使って同じことをする．声の一致度から顔の一致度を引く．
\end{enumerate}

\paragraph{差し引く分}
一致度が高くても，声の特徴が言葉の情報を持っているとは限らない．次の 2 つの理由だけでも一致は起きる．
\begin{itemize}
  \item \textbf{両方が同じ感情を当てた}：本当に Happy の発話なら，言葉の分類器も Happy と答えやすく，声の特徴も Happy の手がかりを持つ．言葉の情報が無くても，両方が正解すれば一致する．
  \item \textbf{同じ話者なので判断が似た}：ある話者の発話を言葉の分類器がいつも Sad と答える癖があり，声の特徴がその話者を見分けられれば，一致する．
\end{itemize}
この分を取り除くため，言葉の分類器の判断を，「本当の感情が同じで，話者（MOSEI はチャンネル）が同じ発話」の中でだけでたらめに入れ替えた．入れ替えた判断にも，感情と話者で決まる分は残るが，発話ごとの判断の違いは消える．元の判断の一致度から入れ替えた判断の一致度を引けば，「感情と話者で説明できない部分を真似できた量」になる．

\paragraph{AFEW を数えから外した理由}
AFEW の言葉は，音声を Whisper で書き起こして作った．声と言葉が同じ音から来ているので，声の特徴が言葉の判断を真似できても，それは書き起こしの作り方による．このため AFEW は判定の数えに入れなかった．

\subsection{結果}

\begin{figure}[H]
\centering
\begin{tikzpicture}
\begin{axis}[
  width=0.9\textwidth, height=5.8cm,
  ybar=3pt, bar width=13pt,
  ymin=20, ymax=58,
  ylabel={一致度（点）},
  symbolic x coords={AFEW,IEMOCAP,MELD,MOSEI},
  xtick=data, enlarge x limits=0.2,
  legend style={at={(0.5,1.03)},anchor=south,legend columns=2,font=\small,draw=none},
  ymajorgrids, grid style={black!10},
  tick label style={font=\small}, label style={font=\small},
  nodes near coords, nodes near coords style={font=\scriptsize},
  every node near coord/.append style={/pgf/number format/fixed,/pgf/number format/fixed zerofill,/pgf/number format/precision=1},
]
\addplot[fill=cA,draw=cA] coordinates {(AFEW,41.4) (IEMOCAP,52.8) (MELD,46.6) (MOSEI,49.0)};
\addplot[fill=cV,draw=cV] coordinates {(AFEW,36.5) (IEMOCAP,36.1) (MELD,31.7) (MOSEI,33.5)};
\draw[black!70,dashed] (rel axis cs:0,0.1316) -- (rel axis cs:1,0.1316) node[pos=0.01,above right,font=\scriptsize]{でたらめ（25 点）};
\legend{声の特徴で真似,顔の特徴で真似}
\end{axis}
\end{tikzpicture}
\caption{言葉の分類器の判断を，声の特徴と顔の特徴で真似したときの一致度（K = 5）．AFEW は言葉が音声の書き起こしなので判定の数えに入れていない．}
\label{fig:mimic}
\end{figure}

\begin{table}[H]
\centering
\caption{言葉の分類器の判断の真似：声の特徴の一致度 − 顔の特徴の一致度（K = 5）．単位は点．}
\label{tab:mimic}
\small\setlength{\tabcolsep}{4pt}
\begin{tabular}{lcccc}
\toprule
 & AFEW（数えの外） & IEMOCAP & MELD & MOSEI \\
\midrule
一致度の差 & +4.9 [\textminus{}2.7, 12.5] & \textbf{+16.7} [7.3, 26.1] & \textbf{+14.9} [12.7, 17.0] & \textbf{+15.5} [14.4, 16.5] \\
感情と話者の分を & +10.7 [1.9, 19.6] & \textbf{+17.3} [8.6, 26.0] & \textbf{+13.9} [11.2, 16.7] & \textbf{+10.2} [9.1, 11.3] \\
\quad 差し引いた後の差 & & & & \\
\bottomrule
\end{tabular}
\end{table}

図\ref{fig:mimic}と表\ref{tab:mimic}に結果を示す．
\begin{itemize}
  \item 声の特徴は，言葉の分類器の判断を，顔の特徴よりよく真似できた．差は IEMOCAP で +16.7 点，MELD で +14.9 点，MOSEI で +15.5 点で，3 つとも区間は 0 を含まない．
  \item 感情と話者の分を差し引いても差は残った．IEMOCAP で +17.3 点，MELD で +13.9 点，MOSEI で +10.2 点である．
  \item 計算の前に決めた基準は「AFEW を除く 3 コーパスのうち 2 つ以上で差が 0 より大きい」で，3 つとも満たした．
\end{itemize}

\subsection{確かめられなかったこと}

研究者の説明には，「声だけが，ほかのモダリティの手がかりを持っている」という部分もあった．これが正しければ，逆向きの「声の特徴は，顔の分類器の判断を，言葉の特徴より良く真似できる」も成り立つはずである．この予言の基準は「4 コーパスのうち 3 つ以上で差が 0 より大きい」と計算の前に決めた．結果は，差し引く前で 4 コーパスのうち 2 つ（MELD と MOSEI），差し引いた後で 1 つ（MELD）しか基準を満たさなかった．このため，「声だけが」の部分は主張から外した．

\subsection{判定には使っていないが，見えたこと}

言葉の特徴も，声の分類器の判断をよく真似できた．感情と話者の分を差し引いた後の一致度は 4 コーパスで 4.7〜14.2 点で，すべて 0 より大きかった．声と言葉はお互いの判断を真似でき，顔だけがこの 2 つから離れている．

%======================================================================
\section{考察}

\subsection{声が言葉の判断を真似できる理由}

声の特徴（WavLM の最後の層）が言葉の分類器の判断を真似できたことは，「音声には話している言葉の情報も入っている」という研究者の説明と合う．音声の自己教師あり学習のモデル（WavLM など）の内部に単語の情報があることは先行研究でも報告されている（Pasad ほか 2023 は，中間の層ほど単語の情報が多いと報告した）．

ただし，「声の特徴に言葉の中身が入っているから真似できた」とまでは言えない．感情と話者の分を差し引いた後に残る一致には，言葉の中身のほかに，次の共通の手がかりも入りうる．
\begin{itemize}
  \item 覚醒度（声の勢い．興奮しているか落ち着いているか）．怒って強い言葉を使う人は，声も大きくなりやすい．
  \item 発話の長さ，録音の場面．
  \item 書き起こしの記号．言葉の文字には「!」や笑い声の注記などが入ることがあり（MELD では 3 モダリティがそろう発話の 26.9\%），これは声の手がかりの一部を言葉に写したものである．ただし，こうした記号の無い発話だけで測り直しても，主張 2 の判定は変わらなかった．
\end{itemize}

\subsection{顔と融合の分類器が上がらない理由}

\begin{itemize}
  \item \textbf{顔}：顔の特徴は，言葉の分類器の判断も声の分類器の判断も真似しにくかった（図\ref{fig:mimic}と 5.5 節）．疑似ラベルは 3 つのモダリティの予測の積なので，顔の分類器から見ると，自分の特徴では説明できない判断を多く含む．AFEW のように顔がもともと最も強いコーパスでは，疑似ラベルを真似ると逆に損をしうる．
  \item \textbf{融合}：融合の分類器は初めから 3 つのモダリティを全部見ているので，疑似ラベルは融合の分類器に新しい情報をほとんど与えない．
\end{itemize}

\subsection{提案と普通の蒸留がほぼ同じ値になる理由}

提案の UAR は，同じ疑似ラベルを使った温度つき蒸留とほぼ同じだった（3 節）．3 つの条件から見ると，学び直しで上がる分は「余地」「疑似ラベルの当たり」「真似できる分」で決まり，疑似ラベルをどの損失の形で渡すかではほとんど変わらない，と解釈できる．この解釈は直接には確かめていない．このため論文は，手法の優劣ではなく，疑似ラベルで学び直す方法の全体について「どの分類器で，なぜ上がるか」を説明する解析として書く．

\subsection{疑似ラベルの誤り方を差し引く試み（実験 15）}

3 つの条件のうち条件 2（疑似ラベルの当たり）を良くできれば，UAR がさらに上がるかもしれない．そこで，K 本（20 本）の上で「疑似ラベルが感情をどう取り違えるか」の表を数え，その取り違えを差し引いてから学び直す方法を試した（既知の方法である Gold Loss Correction の形）．設定は仮のターゲットの回で選んで固定し，評価の回を 1 回だけ流した．
\begin{itemize}
  \item 仮のターゲットの回の選び方は，声の分類器の 4 コーパスすべてで「差し引かない」を選んだ．
  \item 選び方を通さずに，K 本で数えた表で差し引いた版を評価の回で比べると，声の分類器は 4 コーパスのうち 3 つ（AFEW \textminus{}1.3 点，MELD \textminus{}4.8 点，MOSEI \textminus{}2.6 点）で差し引かない版より下がり，上がったコーパスは無かった．
  \item 評価のときだけ使える全部のラベルで数えた正確な表を使っても，声の分類器で上がったのは MOSEI の +1.0 点 [0.2, 1.8] だけだった．言葉の分類器は，正確な表なら AFEW +1.8 点，MELD +2.3 点，MOSEI +3.6 点上がった．
  \item K = 5 では 20 本で表を数えるので表の揺れが大きく，MOSEI では表の 200 行のうち 86 行で，正しい感情の欄が別の感情の欄より小さくなった．
\end{itemize}
この結果から，声の分類器では，疑似ラベルを感情ごとに直して当たりを増やす余地は小さい（正確な表でも 1 点ほど）．さらに，その表を K 本から作ることも難しい．

その後，ソースのコーパスで数えた取り違えの表を使う案も検討した．別の検査役が仮のターゲットの回の表だけで集計すると，表の形はコーパスごとに違い（同じコーパスの中で組み合わせを替えたときの表のずれは平均 0.017，コーパスを替えたときは平均 0.101），ほかのコーパスの表を使っても表は正確にならなかった．この案は実験していない．この段落の数は検査役の集計で，結果の照合はまだ通していない．

\subsection{今後の問い}

条件 3 に近い量として，「疑似ラベルの判断を自分の特徴で真似したときの一致度」（5 節の一致度と同じ作り．差し引く前のもの）は，ターゲットのラベルを使わずにラベルの無い発話だけで測れる．学び直す前にこの量を測って，学び直すかどうかを決められる可能性がある．ただし，評価の回で「真似できる分が大きいほど gain が大きい」という関係を計算の前に決めた規則で確かめたところ，判定は「判定できない」だった（12 の組み合わせのうち，説明どおりの関係が出たのは 5）．この使い方は，別のデータで確かめる必要がある．

%======================================================================
\section{弱いところ（論文の限界に書くこと）}

\begin{itemize}
  \item 3 つの条件は，結果を見た後に立てた説明である．使ったのは同じ 4 コーパスと 1 組のエンコーダ（WavLM，AffectNet で学習した ResNet-50，RoBERTa）だけで，別のデータでは確かめていない．
  \item 方法の設定は，評価に使ったのと同じ 4 コーパスの上で（ソースを仮のターゲットにして）決めた．
  \item 「声が特別」とは書かない．書けるのは「この 4 コーパスとこのエンコーダでは，3 つの条件がそろったのは声の分類器だった」までである．
  \item 疑似ラベルには，学び直す分類器自身の予測も入っている．声の上がり幅を，ほかのモダリティの情報だけに帰すことはできない．
  \item 3 つの条件の数の一部（表\ref{tab:cond}の「欠けた条件」の判断）は下調べの値に基づく．評価の回では声の真似できる分が下調べより小さく，IEMOCAP の声では区間が 0 を含んだ．
  \item 主張 2 の差し引きは，感情と話者（MOSEI はチャンネル）で説明できる分だけを除く．覚醒度，録音の場面などの共通の手がかりは残る．MELD の映像には話者でない人の顔が写っていることがある．IEMOCAP の映像の約半分には動きを測るためのマーカーが写る．
  \item 顔のエンコーダは表情のラベルで学習したものなので，表情以外の手がかりが少ない可能性がある．別の顔のエンコーダ（CLIP など）では確かめていない．
  \item エンコーダを凍結し，線形の分類器だけを学んだ．エンコーダごと学習し直す場合には結果が変わりうる．
\end{itemize}

%======================================================================
\section{まとめ：論文で主張すること}

\begin{yoten}{論文に書ける主張}
\begin{enumerate}
  \item ほかのモダリティを含む疑似ラベルで線形の感情分類器を学び直したとき，UAR が上がるのは「余地がある」「疑似ラベルが今の分類器より当たる」「疑似ラベルの判断を自分の特徴で真似できる」の 3 つがそろう場合である．この 4 コーパスとこのエンコーダでは，4 コーパスすべてで 3 つがそろったのは声の分類器で，声の分類器の UAR は +2.4〜+4.1 点上がった．
  \item 声の特徴（WavLM の最後の層）は，言葉の分類器の判断を，顔の特徴（ResNet-50）より良く真似できる．その差は IEMOCAP，MELD，MOSEI で +14.9〜+16.7 点で，感情と話者で説明できる分を除いても +10.2〜+17.3 点残る．
\end{enumerate}
\end{yoten}

\begin{yoten}{論文に書かないこと}
\begin{itemize}
  \item 「声だけが，ほかのモダリティの手がかりを持っている」（逆向きの予言が基準を満たさなかった）．
  \item 「声の特徴に言葉の中身が入っている」（差し引いた後の一致に，言葉の中身以外の共通の手がかりが入りうる）．
  \item 「提案は既存の蒸留より良い」（48 の条件のうち 40 で差の区間が 0 を含んだ）．
\end{itemize}
\end{yoten}

\appendix
\section{用語の表}
\begin{center}
\small
\begin{tabular}{>{\raggedright}p{3.2cm}p{11.1cm}}
\toprule
用語 & 意味 \\
\midrule
発話 & 1 人が話した一区切り．数秒の動画 \\
モダリティ & 情報の種類．この研究では声・顔・言葉 \\
コーパス & 感情のラベルが付いた発話の集まり \\
ソース／ターゲット & ラベルが全部ある既存のコーパス／感情ごとに K 本しかラベルが無い新しいコーパス \\
K 本 & ターゲットで感情ごとに付いているラベルの本数．この資料では K = 5（合わせて 20 本） \\
LOCO 設定 & 4 コーパスのうち 1 つをターゲット，残り 3 つをソースにする組み方 \\
エンコーダ，特徴 & 発話を数の並び（特徴）に変える学習済みのネットワーク．この研究では中身を変えない（凍結） \\
線形の分類器 & 特徴の数に重みを掛けて足した点数で感情を決める分類器 \\
UAR & 感情ごとの正解率の平均．でたらめなら 25 点 \\
疑似ラベル & 人ではなく分類器が付けた仮の正解．この研究では 3 つの分類器の予測の積 \\
提案 & K 本と疑似ラベルで分類器を学び直す方法 \\
K 本だけ & 比べる相手．出発点から K 本だけで学び直した分類器 \\
gain & 提案の UAR − K 本だけの UAR \\
余地 & 全部のラベルで学んだ分類器の UAR − K 本だけの UAR \\
当たる幅 & 疑似ラベルの UAR − K 本だけの UAR（ラベルの無い発話の上） \\
真似できる分 & 自分の特徴で疑似ラベルを当てるように学んだ分類器の UAR − K 本だけの UAR \\
一致度 & ある分類器の判断を別の特徴で当てたときの，感情ごとの一致の割合の平均．でたらめなら 25 点 \\
95\% 区間 & 50 回の結果のばらつきから作った，本当の値がおそらく入る範囲 \\
仮のターゲットの回 & ソースの 1 つをターゲットに見立てて設定を選ぶ試し．評価のテストの正解を見ない \\
\bottomrule
\end{tabular}
\end{center}

\section{数の出どころ}
\begin{center}
\footnotesize
\begin{tabular}{>{\raggedright}p{4.0cm}>{\raggedright\arraybackslash}p{10.3cm}}
\toprule
この資料の数 & 出どころ（MER-DA のリポジトリの中．結果のファイルは git に入らず，控えは \texttt{DAI/evidence/主張\_2026-11-12/}） \\
\midrule
表\ref{tab:gain}の gain & \texttt{results/dai1/stats/families\_flagged.csv}（族 1，K = 5）．計算はコミット ba98974．照合 \texttt{DAI/checks/DAI-eval/claim\_v3.md} \\
表\ref{tab:gain}の UAR の水準，蒸留の UAR & \texttt{results/dai1/stats/desc/levels.csv}（記述の出力） \\
提案と蒸留の比べ（48 のうち 40） & \texttt{DAI/docs/結果の文の下書き\_2026-11-05.md} の B5（族 3） \\
表\ref{tab:cond} & \texttt{DAI/docs/実験14\_結果の文の下書き\_2026-11-12.md} 4 節（照合 \texttt{DAI/checks/DAI-14/claim\_v7.md}）．余地は \texttt{results/dai3/eval/expl4\_corpus\_restored.csv}（記述の出力）．「欠けた条件」は \texttt{DAI/docs/主張\_2026-11-12.md} 2 節（下調べ \texttt{results/dai11\_probe/req\_vs\_gain.csv}） \\
融合の当たる幅と一致の割合 & \texttt{DAI/docs/結果の文の下書き\_2026-11-05.md} の E1 \\
表\ref{tab:mimic}，図\ref{fig:mimic}，5.4・5.5 節 & \texttt{results/dai14/summary/predictions.csv}，\texttt{descriptive.csv}．計算はコミット fbd6c07．結果の文 \texttt{DAI/docs/実験14\_結果の文の下書き\_2026-11-12.md}（第 7 版） \\
6.4 節（実験 15） & \texttt{DAI/docs/実験15\_結果の文の下書き\_2026-11-13.md}（第 8 版．照合 \texttt{DAI/checks/DAI-15/claim\_v8.md}） \\
6.4 節の最後の段落 & \texttt{research\_notes/改善案\_2026-11-13/review\_idea\_source\_prior.md}（検査役の集計．照合の前） \\
6.5 節 & \texttt{DAI/docs/主張\_2026-11-12.md} 4 節（\texttt{results/dai3/eval/expl3.csv}） \\
\bottomrule
\end{tabular}
\end{center}

\end{document}
